\section{The detector and a finite Euler bound}\label{sec:detector}

This section establishes the detection statement needed for the tail
argument.  Its conclusion concerns actual sphere classes, including
their finite additive order.  We also prove a bound on the length of
the Euler-localized homotopy fixed point spectral sequences.  Neither
argument requires an explicit presentation of the stabilizer action
on the deformation parameters.

\subsection{The action and grading conventions}

Let \(k=\F_{3^6}\), and let \(E=E_6\) be Morava \(E\)-theory for the
Honda formal group of height six over \(k\), with its coherent
stabilizer action.  We use the standard Lubin--Tate and
Goerss--Hopkins--Miller constructions; see \cite{LubinTate,GH}.
Write
\begin{equation}\label{det:coefficients}
 E_* = R[u^{\pm1}],\qquad
 R=W(k)[[u_1,\ldots,u_5]],\qquad
 |u|=-2,\qquad
 \m=(3,u_1,\ldots,u_5).
\end{equation}
The cotangent line of the universal formal group is \(E_2\):
it is obtained by restricting the augmentation ideal of
\(E^0(\mathbf{CP}^{\infty})\) to \(\mathbf{CP}^1\).
Its dual, the Lie line, is therefore \(E_{-2}\).
The choice of \(u\) trivializes this line nonequivariantly.

The height classification and stabilizer description originate in
the work of Lazard, Dieudonn\'e, Honda, and Lubin;
see \cite[Theorem~IV]{Lazard} and \cite[Theorem~1]{Honda}.
We use the finite-field presentation in
\cite[Appendix~A2.2.12 and A2.2.17--18]{Ravenel}.
Let \(D\) be the central division algebra over \(\Q_3\) of invariant
\(1/6\), with maximal order
\[
 \calO_D=
 W(k)\langle S\rangle/(S^6=3,\;Sa=a^\varphi S),
\]
where \(\varphi\) is Witt Frobenius.  Passing to the opposite algebra,
if required by the action convention, changes none of the arguments
below.  Normalize its valuation by \(v_S(S)=1\).
The degree-six extension \(\Q_3(\zeta_9)\) splits \(D\) and embeds as a
maximal subfield.  Its ring of integers embeds in \(\calO_D\).
For the splitting criterion, centralizer dimensions, and
Skolem--Noether conjugacy used here and below, see
\cite[Chapter~IV, Theorems~2.10 and~3.1, Corollary~3.7, and Remark~4.4(c)]{MilneCFT}.
We fix the resulting subgroups
\begin{equation}\label{det:groups}
 G=\langle g\rangle\cong C_9,\qquad
 H=\langle \sigma\rangle\cong C_3,\qquad
 \sigma=g^3 .
\end{equation}
The cyclotomic uniformizer gives
\begin{equation}\label{det:valuations}
 v_S(g-1)=1,\qquad v_S(\sigma-1)=3.
\end{equation}
For example, if \(\pi=g-1\), then
\(\sigma-1=3\pi+3\pi^2+\pi^3\), whose last term has strictly
smallest \(S\)-valuation.  These are subgroups of the nonextended
stabilizer: their actions on \(R\) fix \(W(k)\).

Choose a homotopy multiplicative, unital \(3\)-typical complex
orientation, normalized to the Honda coordinate on the special
fibre.  For \(1\leq i\leq5\), the elements
\[
 u_i=v_i u^{3^i-1}
\]
give deformation parameters.  We use Hazewinkel generators when an
integral cobar calculation is required.  With right-unit conventions,
the first strict coordinate-change coefficients satisfy
\begin{equation}\label{det:right-units}
 \eta_R(v_1)=v_1+3t_1,\qquad
 \eta_R(v_2)\equiv
 v_2+v_1t_1^3-v_1^3t_1\pmod3 .
\end{equation}
In particular, \(t_1\), of internal degree \(4\), is a
one-cocycle.  Evaluation of the cobar complex on the action uses these
strict changes of the chosen graded coordinate.

All cohomological spectral sequences have bidegrees \((s,t)\), with
\[
 d_r:E_r^{s,t}\longrightarrow E_r^{s+r,t+r-1}.
\]
Thus \(t-s\) is the homotopy degree.  Representation suspensions will
be specified separately below.

\subsection{Comparison of filtrations}

Write
\[
 \mathcal N_2^{s,t}
 =\Ext_{\BP_*\BP}^{s,t}(\BP_*,\BP_*)
\]
for the ordinary \(3\)-local Adams--Novikov \(E_2\)-term.
The Thom reduction is denoted
\[
 \Phi:\mathcal N_2^{s,t}
 \longrightarrow
 \Ext_{\calA_3}^{s,t}(\F_3,\F_3).
\]

\begin{lemma}[Comparison]\label{det:filtration-comparison}
There are augmented comparisons from the ordinary
Adams--Novikov spectral sequence to the mod-\(3\) Adams spectral
sequence and to the homotopy fixed point spectral sequences for
\(E^{hG}\) and \(E^{hH}\).
If a sphere class \(\alpha\) is detected by \(b_j\), then
\(\alpha\) has Adams--Novikov filtration exactly two.
It admits a permanent representative
\[
 x\in\mathcal N_2^{2,\,4\cdot3^{j+1}},
 \qquad \Phi(x)\in\F_3^\times b_j .
\]
\end{lemma}

\begin{proof}
The unit-preserving map \(\BP\to\HF\) induces a map between the
unit fibres and their smash powers.  Consequently the filtrations
on sphere homotopy satisfy
\begin{equation}\label{det:filtration-inclusion}
 F_{\mathrm{AN}}^s\pi_*S
 \subseteq F_{\mathrm{Adams}}^s\pi_*S .
\end{equation}
In particular, a class of Adams filtration exactly two has
Adams--Novikov filtration at most two.  The connective
Adams--Novikov spectral sequence converges to \(3\)-local sphere homotopy
\cite[Theorems~2.2.3 and~2.2.13--14]{Ravenel}; see also
\cite[p.~475, after Equation~(1.10)]{MRW}.
These positive stable stems are finite, so their \(3\)-local and
\(3\)-complete groups agree.  Its \(E_2\)-term
has internal degree divisible by \(4\).  In the stem
\(4\cdot3^{j+1}-2\), filtration zero or one is therefore impossible.
Thus the filtration is exactly two.
The induced map on \(E_2\) is Thom reduction.
There is no incoming Adams differential in the bidegree of \(b_j\):
the potential length-two source has filtration zero and positive
odd internal degree, and longer sources have negative filtration.
The associated-graded image of the chosen class consequently
gives the asserted nonzero multiple of \(b_j\).

For completeness, the comparison with group cochains can be
constructed without requiring an equivariant chosen orientation.
Use the semicosimplicial Amitsur object with level
\(\MU_{(3)}^{\wedge(s+1)}\).
For \(K=G\) or \(H\), apply the translated orientation to the
\(\ell\)-th factor at \(g_\ell\), and multiply in \(E\):
\begin{equation}\label{det:homogeneous-comparison}
 (a_0,\ldots,a_s)\longmapsto
 \bigl[(g_0,\ldots,g_s)\longmapsto
       \prod_{\ell=0}^s g_\ell f(a_\ell)\bigr].
\end{equation}
The target is the spectrum of equivariant homogeneous group
cochains.  Deleting a group entry corresponds to inserting the unit
in the indicated source factor.
The orientation is a map under the unit; its translates form a
coherent family under that unit.  Their exterior powers, followed
by the coherent commutative multiplication in \(E\), thus make
\eqref{det:homogeneous-comparison} compatible with all cofaces and
with the augmentation from the sphere.
Only unit insertions occur here, so no \(E_\infty\) refinement of
the orientation is required.

Forgetting degeneracies gives the unnormalized rather than the
normalized cochain complex; the degenerate subcomplex is
contractible.  It therefore leaves the comparison unchanged from
\(E_2\) onward.  The finite-stage comparison also respects actual
filtration-two lifts.  The ordinary totalization tower maps under
the augmentation to its fat tower.  In the target,
\(s^0:C^1\to C^0\) has section \(d^0\), and \(d^1-d^0\)
lands in \(N^1=\ker s^0\).  Hence
\[
 \Tot^{\mathrm{fat}}_1 C^\bullet
 \simeq \Tot_1 C^\bullet\oplus\Omega C^0,
\]
with an augmentation-compatible projection onto the first summand.
Composing the fat group-cochain comparison with this projection
therefore gives an actual ordinary filtration-two lift, whose
\(H^2\)-symbol is the normalized comparison class.
Quillen's \(3\)-typical splitting
\cite[Theorems~4--5]{Quillen} identifies the
\(\MU_{(3)}\) and \(\BP\) descriptions of the ordinary
Adams--Novikov spectral sequence.
On the \(\BP\) cobar, the comparison is precisely substitution of
the chosen coordinates and their strict coordinate changes.
The augmentations ensure that it sends actual filtered sphere
classes to the corresponding filtered homotopy fixed point classes.
\end{proof}

\subsection{The Adams--Novikov two-line in the detector degree}

Let \(\delta_3\) denote the integral connecting homomorphism for
multiplication by \(3\), and let \(\delta_{v_1^d}\) denote the
connecting homomorphism over \(\BP_*/3\) for multiplication by
\(v_1^d\).  In the following lemma the prime indicates an
uncorrected representative; it does not assert an identification
with every conventional choice of a Greek-letter generator.

\begin{lemma}\label{det:two-line-basis}
Fix \(j\geq0\).  For
\[
 n\in\{j,j-2,\ldots,j\bmod2\},\qquad
 d=3^n,\qquad L=\frac{3^{j+1-n}+1}{4},
\]
the classes
\begin{equation}\label{det:uncorrected-beta}
 \beta'_{Ld/d}=\delta_3\delta_{v_1^d}(v_2^{Ld})
\end{equation}
form an \(\F_3\)-basis of
\(\mathcal N_2^{2,\,4\cdot3^{j+1}}\).
Every member with \(L>1\) has zero Thom image.
\end{lemma}

\begin{proof}
The two-line theorem of Miller--Ravenel--Wilson
\cite[Theorem~2.6]{MRW}, also stated in
\cite[Corollary~5.4.5(a)]{Ravenel}, describes a direct sum of cyclic
groups with parameters \(m=L3^n\), \(3\nmid L\), denominator
\(v_1^d\), and order \(3^{i+1}\).  The internal degree is
\(4(4m-d)\).  Its restrictions include
\[
 0\leq i\leq n,\qquad
 1\leq d\leq a_{n-i},\qquad
 a_0=1,\qquad a_r=3^r+3^{r-1}-1\quad(r\geq1).
\]
In particular \(d<2\cdot3^n\).
In the desired degree the equation is
\[
 4L3^n-d=3^{j+1}.
\]
If \(n\geq j+1\), it gives \(d\geq3\cdot3^n\), a contradiction.
For \(n\leq j\), division by \(3^n\) shows that the positive
integer \(4L-3^{j+1-n}\) is less than \(2\), hence is \(1\).
Thus \(d=3^n\), \(L\) has the displayed value, and integrality
forces \(n\equiv j\pmod2\).
For \(n=0\), the bound \(i\le n\) gives \(i=0\); for \(n\ge1\),
the inequality \(a_{n-1}<3^n\) again forces \(i=0\).
This bounds the dimension by the number of classes in
\eqref{det:uncorrected-beta}, and shows that the group has
exponent \(3\).
The remaining index restrictions are satisfied by all these candidates:
\(3^i\mid d\) is automatic, the special bound \(d\leq3^n\) when
\(L=1\) is an equality, and maximality of \(i\) follows from
\(a_{n-1}<3^n\) for \(n\geq1\), while \(3\nmid d\) when \(n=0\).

We next prove independence of these uncorrected classes.
Equation~\eqref{det:right-units} and the Frobenius identity in
characteristic \(3\) show that \(v_2^{Ld}\) is primitive modulo
\((3,v_1^d)\), so its inner connecting class
\[
 c_n=\delta_{v_1^d}(v_2^{Ld})
 \in\Ext_{\BP_*\BP}^{1,\,4\cdot3^{j+1}}(\BP_*,\BP_*/3)
\]
is defined.
Its exact \(v_1\)-annihilator exponent is \(d\).
Indeed \(v_1^d c_n=0\), whereas the connecting-sequence comparison
gives
\[
 v_1^{d-1}c_n=\delta_{v_1}(v_2^{Ld})\ne0.
\]
For the last assertion, the primitives of \(\BP_*/3\) are
\(\F_3[v_1]\).  One can see this directly: if a primitive
polynomial depends on a highest variable \(v_s\), \(s\geq2\),
set all positive \(t_i\) except \(t_{s-1}\) to zero.
The lower \(v_i\) are fixed, while the right unit sends
\[
 v_s\longmapsto
 v_s+v_1t_{s-1}^3-v_1^{3^{s-1}}t_{s-1}.
\]
Comparing the highest power of \(t_{s-1}\) in a polynomial
depending nontrivially on \(v_s\) contradicts invariance.
The image of the primitives in \(\BP_*/(3,v_1)\) is therefore
constant, and does not contain \(v_2^{Ld}\).
This proves the nonzero connecting class.
Since the \(d\)'s are distinct, multiplication by
\(v_1^{d_{\max}-1}\) in a proposed relation isolates its
largest-exponent summand.  Hence the \(c_n\)'s are independent.

It remains to exclude relations created by \(\delta_3\).
Set \(N=3^{j+1}\).  The integral alpha-line calculation
\cite[Theorem~2.2]{MRW} says that the image of integral
\(\Ext^1\) in this internal degree is spanned, after reduction
modulo \(3\), by
\begin{equation}\label{det:alpha-reduction}
 v_1^{N-1}t_1 .
\end{equation}
Explicitly, an integral representative is the differential of
\(v_1^N\) divided by \(3^{1+\nu_3(N)}\).
Expanding \((v_1+3t_1)^N-v_1^N\), its linear term reduces to
\eqref{det:alpha-reduction}; its \(k\)-th term for \(k\geq2\)
has \(3\)-adic valuation
\(k-1-\nu_3(k)>0\).
The class \eqref{det:alpha-reduction} is not \(v_1\)-torsion.
Evaluate it on the height-one Honda formal group in characteristic
\(3\) at the strict automorphism \(1+3\).
The coefficient \(v_1\) is then invertible, the coefficient action
is trivial, and the first strict coefficient is nonzero.
Evaluation of a one-coboundary at that automorphism is zero,
whereas every \(v_1\)-multiple of
\eqref{det:alpha-reduction} evaluates nontrivially.
The kernel of \(\delta_3\) consequently meets the
\(v_1\)-torsion span of the \(c_n\)'s trivially.
Their images are independent and attain the previously established
dimension bound.

To compute the Thom image, put
\(I=(3,v_1,v_2,\ldots)\) in the left coefficient ring of the
cobar complex.  Modulo \(3\), write
\[
 A_1=t_1^3-v_1^2t_1,\qquad
 c_n=\frac{(v_2^d+v_1^d A_1^d)^L-v_2^{Ld}}{v_1^d}.
\]
Every term of this polynomial has at least \((L-1)d\)
positive coefficient factors.  It therefore admits an integral
lift in \(I^{(L-1)d}\).
The ideal \(I\) is invariant, so the cobar differential preserves
its powers.  Moreover, on each free left coefficient module,
\[
 (I^q:3)=I^{q-1}\qquad(q\geq1);
\]
this follows, for example, because the initial form of \(3\)
is a non-zero-divisor in the \(I\)-graded polynomial ring.
The integral differential of the lift is divisible by \(3\);
after division it still lies in \(I^{(L-1)d-1}\).
When \(L>1\), we have \((L-1)d\geq2\), so reduction of
this representative under Thom reduction is zero.
\end{proof}

\begin{remark}\label{det:corrected-beta}
The degree calculation uses the restrictions in the full
two-line theorem, while the independence argument provides the
specified uncorrected basis.  Thus correction terms in the
Miller--Ravenel--Wilson representatives have not been suppressed.
For example, the two basis directions in the degree of \(b_2\)
are the uncorrected classes \(\beta'_{9/9}\) and \(\beta'_7\);
a permanent combination need not be the conventional
\(\beta_{9/9}\) alone.
\end{remark}

\subsection{Integral CM detection and the actual image}

\begin{lemma}[Integral CM evaluation]\label{det:cm-evaluation}
For \(j\geq2\), every
\[
 x\in\mathcal N_2^{2,\,4\cdot3^{j+1}}
 \quad\text{with}\quad
 \Phi(x)\in\F_3^\times b_j
\]
has nonzero image in \(H^2(G;E_{4\cdot3^{j+1}})\).
\end{lemma}

\begin{proof}
We use a fixed Lubin--Tate deformation to evaluate the
comparison of Lemma~\ref{det:filtration-comparison}.
Let \(\zeta\) be a primitive ninth root of unity, and let
\(\calO=W(k)[\zeta]\), allowing a residue extension if necessary.
The Lubin--Tate construction for \(\Z_3[\zeta]\)
\cite[Theorem~1]{LubinTateLocal}
gives a height-one formal \(\Z_3[\zeta]\)-module over
\(\calO\).  Its underlying special fibre has height six.
After residue extension it can be identified with the chosen
Honda formal group.  The embedding of
\(\Q_3(\zeta)\) in its rational endomorphism algebra is conjugate
to our chosen embedding by Skolem--Noether.
The conjugator can be made a unit: multiply it by an element
of the centralizing cyclotomic field whose \(S\)-valuation
cancels its own, using \(v_S(\zeta-1)=1\).
This gives a framing identifying the ninth-root action with \(G\).

We obtain a \(G\)-equivariant evaluation of graded coefficient
rings
\[
 E_*\longrightarrow \calO[w^{\pm1}],\qquad |w|=-2.
\]
The group fixes \(\calO\), and acts on \(w\) by \(\zeta\)
or its inverse, according to the action convention.
The elements \(v_1,v_2\) have zero residue modulo
\(\pi=\zeta-1\).
By \eqref{det:valuations}, the first strict coefficient
\(t_1(g)\) has nonzero residue.
The latter statement concerns the strict coordinate change
after grading: its residual automorphism has linear coefficient
\(1\), and \(g-1\) first occurs at the third power of the
coordinate.

Put \(T=4\cdot3^{j+1}\).
The degree-\(T\) line has character
\(\zeta^{\pm2\cdot3^{j+1}}=1\), hence trivial \(G\)-action.
After trivializing this line, the periodic cyclic cochain
complex is
\[
 \calO\xrightarrow{0}\calO\xrightarrow{9}\calO
       \xrightarrow{0}\calO\xrightarrow{9}\cdots.
\]
Its reduction modulo \(3\) has zero differentials.
The connecting homomorphism for multiplication by \(3\)
is therefore
\begin{equation}\label{det:integral-bockstein}
 \delta_3:H^1(G;\calO/3)\longrightarrow H^2(G;\calO),
 \qquad
 \bar c\longmapsto3c\pmod9 .
\end{equation}
This formula follows by lifting the one-cochain \(c\),
taking its differential \(9c\), and dividing by \(3\).
The inner connecting homomorphism \(\delta_{v_1^d}\) is
formed in the \(\BP\) cobar complex before evaluation.
We evaluate its already formed polynomial one-cocycle;
we do not form a connecting homomorphism for multiplication
by the evaluated \(v_1^d\) on \(\calO/3\).
The comparison from the bar resolution to the cyclic
resolution evaluates a one-cocycle at \(g\), so it computes
the image of the inner connecting class in
\eqref{det:uncorrected-beta}.

Express \(x\) in the basis of Lemma~\ref{det:two-line-basis}.
Every summand with \(L>1\) has zero Thom image; therefore the
coefficient of the \(n=j,L=1\) summand is nonzero.
The residue of the inner cocycle at \(g\) is
\[
 L\,v_2^{(L-1)d}t_1(g)^{3d}.
\]
For \(L>1\) it is zero, whereas for \(L=1\) it is a unit.
Thus the inner cocycle of the entire combination evaluates
to an element \(c\in\calO/3\) with unit residue.
Equation~\eqref{det:integral-bockstein} sends it to
\(3c\ne0\) in \(\calO/9\).
The integral coefficient modules and their cochain complexes
are \(3\)-torsion-free, so evaluation commutes with the
outer connecting homomorphism \(\delta_3\).
This naturality proves nonvanishing in \(H^2(G;E_T)\).
\end{proof}

\begin{remark}\label{det:integral-target}
This is the CM detection mechanism used in
\cite[Section~6]{BR}, with the integral target retained.
The further reduction
\(\calO/9\to\calO/3\) kills \(3c\).
Consequently one must test nonvanishing in
\(H^2(G;\calO)\) before making that reduction.
No coefficient-action presentation is used in
Lemma~\ref{det:cm-evaluation}.
\end{remark}

\begin{lemma}[A filtration-two symbol detects its lift]
\label{det:two-layer}
Let a finite group \(J\) act on an even spectrum \(Z\), and put
\(Y=Z^{hJ}\).
Let \(F^s=\operatorname{fib}(Y\to\Tot_{s-1}C^\bullet)\)
be the decreasing tower of its group-cochain resolution.
If a class in \(\pi_{T-2}F^2\), with \(T\) even, has nonzero
symbol in \(H^2(J;\pi_T Z)\), its image in \(\pi_{T-2}Y\)
is nonzero.
\end{lemma}

\begin{proof}
If the image were zero, the class would lie in the image of
\(\pi_{T-1}(Y/F^2)\), and \(Y/F^2\simeq\Tot_1C^\bullet\).
In the fibre sequence
\[
 \Omega N^1C^\bullet\longrightarrow
 \Tot_1C^\bullet\longrightarrow C^0,
\]
the degree-\((T-1)\) group of \(C^0\) vanishes by evenness.
Every such element therefore comes from
\(\pi_T N^1C^\bullet\).
Its boundary in the second layer is the cochain differential
from degree one to degree two.  The given symbol would thus
be a \(d_1\)-boundary, contrary to its nonzero class in
\(H^2\).
Equivalently, after \(E_2\) the only possible ordinary incoming
differential has length two and starts in odd internal degree;
larger incoming lengths start in negative filtration.
The finite two-layer argument proves the assertion directly,
without a convergence assertion for any localized tower.
\end{proof}

\begin{theorem}[Detection of actual sphere classes]
\label{det:actual-detector}
Let \(j\geq2\), and let
\(\alpha\in\pi_{4\cdot3^{j+1}-2}S_{(3)}\) be detected by
\(b_j\) in the mod-\(3\) Adams spectral sequence.
Its image under the unit
\[
 \pi_{4\cdot3^{j+1}-2}S_{(3)}
 \longrightarrow
 \pi_{4\cdot3^{j+1}-2}(E^{hG})
\]
is nonzero and has finite additive order.
Its filtration-two symbol is the nonzero image of the
Adams--Novikov representative in
Lemma~\ref{det:filtration-comparison}.
\end{theorem}

\begin{proof}
Choose the actual Adams--Novikov filtration-two lift of
\(\alpha\) given by Lemma~\ref{det:filtration-comparison}.
The augmented comparison sends it to a class in
\(\pi_{4\cdot3^{j+1}-2}F^2\).
Its symbol is nonzero by Lemma~\ref{det:cm-evaluation}.
Lemma~\ref{det:two-layer} proves that its image in
\(E^{hG}\) is nonzero.
Finally, a positive stable sphere stem is finite, and a
homomorphism preserves finite additive order.
\end{proof}

\subsection{The full-gap assertion and its torsion-free repair}

The detection theorem produces a nonzero finite-order image from a
sphere class detected by \(b_j\), \(j\geq2\).  The proposed full gap
in degree minus two cannot hold: we now construct infinite-order
classes in every even degree.  Under the proposed \(972\)-periodicity,
however, torsion-freeness in degree minus two would still exclude the
higher detectors.  We prove the full-gap failure using independence
of field automorphisms and the additive transfer, retaining the actual
homotopy class supplied by the transfer rather than only an invariant
coefficient.

\begin{lemma}[Semilinear trace]\label{det:semilinear-trace}
Let a finite group \(J\) act on a characteristic-zero domain \(A\)
and semilinearly on a free rank-one module \(Ae\).
Suppose the kernel of \(J\to\Aut(A)\) acts trivially on \(Ae\).
Then the additive trace
\[
 \Tr_J:Ae\longrightarrow (Ae)^J,\qquad
 z\longmapsto\sum_{\gamma\in J}\gamma z
\]
is nonzero.
\end{lemma}

\begin{proof}
Put \(F=\operatorname{Frac}(A)\), let
\(\alpha_1,\ldots,\alpha_m\) be the distinct automorphisms of \(F\)
induced by \(J\), and choose representatives \(\gamma_i\).
Write \(\gamma_i(e)=c_i e\), where \(c_i\in A^\times\), and let
\(h\) be the order of the kernel.  On \(Fe\) the trace is
\[
 fe\longmapsto
 h\sum_{i=1}^m c_i\alpha_i(f)e .
\]
Distinct field automorphisms are linearly independent over \(F\).
Indeed, from a nonzero relation with the fewest terms, evaluating at
\(af\) and subtracting \(\alpha_m(a)\) times the relation at \(f\)
removes its last term.  Choosing \(a\) on which another
\(\alpha_i\) differs from \(\alpha_m\) gives a shorter nonzero
relation, a contradiction.
Since \(hc_i\ne0\), the displayed trace is nonzero.

Choose \(f\in F\) on which it is nonzero, and write \(f=a/b\).
The nonzero invariant
\(\prod_{\gamma\in J}\gamma(b)\) clears the denominator of \(f\).
Multiplication by this invariant commutes with the trace, and
therefore gives a nonzero trace of an element of \(Ae\).
\end{proof}

\begin{lemma}[Homotopy transfer]\label{det:homotopy-transfer}
Let a finite group \(J\) act coherently on a spectrum \(Z\).
For every integer \(d\), the additive trace on \(\pi_d Z\)
factors as
\[
 \pi_d Z\longrightarrow \pi_d(Z^{hJ})
       \longrightarrow \pi_d Z .
\]
\end{lemma}

\begin{proof}
Work in the category of spectra with coherent \(J\)-action.
Induction and coinduction from the trivial subgroup agree: their
underlying objects are the finite coproduct and finite product of
\(|J|\) copies of a spectrum, which agree naturally in a stable
category.  A map \(f:S^d\to Z\) therefore gives a composite
\[
 \operatorname{triv}(S^d)
 \xrightarrow{\mathrm{diag}}
 \operatorname{coInd}_e^J S^d
 \simeq \operatorname{Ind}_e^J S^d
 \longrightarrow Z.
\]
The last map is adjoint to \(f\); on the summand indexed by
\(\gamma\), it is \(\gamma f\).
The underlying composite is consequently
\(\sum_{\gamma\in J}\gamma f\).
Maps from a trivial-action sphere to \(Z\) are maps from that sphere
to \(Z^{hJ}\), proving the assertion.  For a finite group the
continuous cochain construction gives the same homotopy limit.
\end{proof}

\begin{proposition}\label{det:full-gap}
For every \(n\in\Z\), the group
\(\pi_{-2n}(E^{hG})\) contains an element of infinite additive order.
In particular, \(\pi_{-2}(E^{hG})\ne0\).
\end{proposition}

\begin{proof}
Each \(Ru^n\) is a semilinear \(R\)-line.  We check the kernel
hypothesis in Lemma~\ref{det:semilinear-trace}.
There is no nontrivial ninth root of unity in \(W(k)\).
To see this first for a cube root, its residue is \(1\), so write it
as \(1+3c\).  The equation
\[
 (1+3c)^3-1=9c(1+3c+3c^2)=0
\]
forces \(c=0\).  Applying this twice handles ninth roots.
If \(q^9=1\) in the power-series domain \(R\), its constant term is
therefore \(1\), and
\[
 (q-1)(1+q+\cdots+q^8)=0
\]
forces \(q=1\), since the second factor has nonzero constant term
\(9\).

Now suppose \(\gamma\in G\) acts trivially on \(R\), and write
\(\gamma(u)=q_\gamma u\).
The group relation gives \(q_\gamma^9=1\), hence \(q_\gamma=1\).
It acts trivially on every \(Ru^n\).
Lemma~\ref{det:semilinear-trace} gives a nonzero integral trace in
\(Ru^n\), and Lemma~\ref{det:homotopy-transfer} lifts it to
\(\pi_{-2n}(E^{hG})\).
Its restriction has infinite order because \(R\) is a
characteristic-zero domain.  The lifted class has infinite order
as well.
\end{proof}

This proposition uses additive traces, which vanish in
\(\widehat H^0(G;E_*)\).  It asserts neither torsion-freeness of a
homotopy group nor periodicity.

\begin{corollary}[The torsion-free repair]\label{det:conditional-repair}
Suppose that the graded groups \(\pi_*(E^{hG})\) are
\(972\)-periodic and that \(\pi_{-2}(E^{hG})\) is
torsion-free.  Then \(b_j\) does not survive for any
\(j\geq4\).
\end{corollary}

\begin{proof}
For \(j\geq4\),
\[
 4\cdot3^{j+1}-2=972\cdot3^{j-4}-2.
\]
A survivor would give a nonzero finite-order element of the
group in this degree by Theorem~\ref{det:actual-detector}.
Periodicity would identify it with a nonzero torsion element
in degree \(-2\), a contradiction.
\end{proof}

The hypotheses of Corollary~\ref{det:conditional-repair} are not
assumed in the sequel.  Instead we determine enough outgoing
differentials to exclude the detected classes directly.

\subsection{Euler localization and finite-page conventions}

For \(K=G,H\), let \(\lambda\) be its faithful complex
one-dimensional representation, compatible under restriction.
We also write \(\lambda\) for its underlying real plane.
Use the cofree genuine \(K\)-spectrum associated to the given
coherent action on \(E\).  Representation-suspended group
cochains give an \(RO(K)\)-graded homotopy fixed point spectral
sequence.  If \(V\in RO(K)\), its total homotopy degree is
\(V-s\), and a differential changes \((s,V)\) by
\((r,r-1)\).

The orientation of a complex representation sphere has trivial
determinant character.  It supplies an invertible symbol
\[
 u_\lambda,\qquad |u_\lambda|=(0,2-\lambda),
\]
on \(E_2\).  This is an orientation symbol in the layers,
not a claimed permanent homotopy orientation.
Let \(b\) denote cyclic cohomological periodicity, with
\(|b|=(2,0)\).  The Euler symbol is
\begin{equation}\label{det:euler-symbol}
 a_\lambda=b\,u_\lambda,\qquad
 |a_\lambda|=(2,2-\lambda).
\end{equation}
It represents the actual representation Euler map, of total
degree \(-\lambda\).

Inverting \(a_\lambda\) makes the \(E_2\)-groups the ordinary
Tate cohomology groups, with the additional orientation
variables.  We call the resulting spectral sequence the
\emph{Tate pages}.  Only their finite pages will be used.
The Euler symbol has no outgoing differential as long as it
remains present: the actual Euler map has a filtration-two
lift, obtained from a nonvanishing section of the associated
line bundle on the \(1\)-skeleton.  Its first obstruction
cocycle is exactly \eqref{det:euler-symbol}.
Localization therefore commutes inductively with differentials
and passage to the next page; if an Euler power is hit, the
localized next page is zero.

Restriction and additive transfer are compatible with these
pages.  To see this without choosing incompatible filtrations,
use the same free \(G\)-space \(EG\) and its skeletons for both
groups.  Restriction is then levelwise, and transfer is the
finite sum over cosets.  The cochain maps satisfy the projection
formula and commute with the differentials.  Cellular
comparisons with other free models induce the usual
isomorphisms from \(E_2\).  Since the faithful Euler classes
restrict compatibly, the same statements hold after
localization.

Only lengths \(r\equiv1\pmod4\) can occur in the computations
below.  Evenness first excludes even \(r\).
The commuting central automorphism \(-1\) of the formal group
acts trivially on \(R\), since it lifts to every deformation,
and by \(-1\) on its Lie line.  Thus it acts on internal degree
\(t=\dim_{\mathbf R}V\) by \((-1)^{t/2}\), while it fixes the
orientation symbols, whose representation shifts have dimension zero.
Equivariance of \(d_r\) then requires
\((-1)^{(r-1)/2}=1\), giving the stated congruence.

\begin{lemma}[A null lift gives a bounded incoming differential]
\label{det:finite-boundary}
Let \(F^{s+1}\to F^s\), \(s\geq0\), be a decreasing tower of
spectra, or of genuine spectra for a finite group.
If \(\alpha\in\pi_VF^s\) maps to zero in \(\pi_VF^0\),
its layer symbol is zero by \(E_{s+1}\).
If that symbol is initially nonzero, its first disappearance
is an incoming differential of length at most \(s\).
\end{lemma}

\begin{proof}
Put \(C^q=F^q/F^{q+1}\), and let \(i,j,k\) be the exact-couple
maps obtained from the fibre sequences for these layers.
The cumulative boundary group on \(E_r\) in filtration \(s\)
is
\[
 B_r^s
 =j_s\ker\bigl(\pi_VF^s\longrightarrow
                       \pi_VF^{s-r+1}\bigr),
 \qquad 1\leq r\leq s+1.
\]
This is the usual exact-couple formula \(B_r=j(\ker i^{r-1})\):
for \(r=1\) it is zero, and deriving the exact couple once
adds the images of elements killed by the next iterate of
\(i\).  Induction gives the displayed expression.
The symbol \(j_s\alpha\) has zero image under \(k\), so it
has no outgoing differential on any page where it is
nonzero.  The hypothesis puts it in \(B_{s+1}^s\), proving
the assertion.
Equivalently, \(\alpha\) comes from
\(\pi_{V+1}(F^0/F^s)\); this finite quotient involves exactly
the intervening \(s\) layers.
\end{proof}

\begin{proposition}[Finite Euler bound]\label{det:euler-bound}
In the cofree genuine spectra associated to \(E\), the actual
faithful Euler classes satisfy
\[
 a_{\lambda|H}^{79}=0,\qquad a_\lambda^{235}=0.
\]
Consequently, the \(H\)-Tate pages kill the unit by a
differential of length at most \(157\), and the \(G\)-Tate
pages kill it by a differential of length at most \(469\).
\end{proposition}

\begin{proof}
We give the construction for \(K=H,G\) simultaneously, retaining
\(H=\langle\sigma\rangle\) and writing \(Q=K/H\).
Let
\[
 R_N=N_e^K\MU,\qquad \rho_K=\text{the real regular representation of }K .
\]
We use the derived genuine norm and its commutative-ring
counit; their indexed multiplication and geometric fixed
point formulas are recalled in \cite[Section~2 and
Appendix~B]{HHR2016}.
Nonequivariantly, \(R_N\) is a smash of \(|K|\) copies of
\(\MU\), labelled by \(K\).
Let
\[
 c\in\pi_{52}\operatorname{Res}_eR_N
\]
be the coefficient of \(x^{27}\) in the strict change between
the orientations labelled \(e\) and \(\sigma\).
Norming this class and then applying the counit gives
\begin{equation}\label{det:norm-class}
 z\in\pi_{52\rho_K}^K R_N.
\end{equation}

\emph{Geometric fixed points.}
We claim that \(\Phi^H z=0\) as a genuine \(Q\)-class.
The geometric fixed points of a free indexed norm descend
to the norm indexed by the orbit set, with its residual group action;
see \cite[Corollary~2.11]{MSZNorms}.
After applying \(\Phi^H\), the factors in the domain of
\[
 N_e^K\operatorname{Res}_eR_N\longrightarrow R_N
\]
are labelled by \(Q\times K\), the target factors by \(Q\),
and the multiplication folds them by
\[
 Q\times K\longrightarrow Q,\qquad (aH,c)\longmapsto acH .
\]
This follows equally by viewing the counit on each free
norm factor as right translation of its labels.
The displayed map factors, equivariantly for \(Q\), as
\[
 Q\times K\longrightarrow Q\times Q\longrightarrow Q.
\]
The first stage is the \(Q\)-norm of the multiplication
\(\MU^{\wedge K}\to\MU^{\wedge Q}\) that identifies labels
in each \(H\)-coset.  Under that multiplication the \(e\)
and \(\sigma\) orientations coincide.
Their strict coordinate change is the identity, so \(c\)
maps to zero before the remaining \(Q\)-norm.
This proves the claim with the residual equivariance retained.

\emph{An Euler multiple vanishes in the norm spectrum.}
Every nontrivial subgroup of \(K\) contains \(H\).
The family of subgroups not containing \(H\) is therefore
just the trivial subgroup.  Isotropy separation gives
\[
 EK_+\wedge R_N\longrightarrow R_N
             \longrightarrow \widetilde{EK}\wedge R_N .
\]
For a spectrum concentrated over \(H\), geometric \(H\)-fixed
points identify its maps from \(S^V\) with genuine \(Q\)-maps
from \(S^{V^H}\); see \cite[Section~2.4]{GreenleesShipley}.
Apply this also after smashing with \(S^{-52\rho_K}\).
The obstruction to lifting \eqref{det:norm-class} through
the first term is precisely its image
\[
 \Phi^H z\in
 \pi_{(52\rho_K)^H}^{Q}\Phi^H R_N,
\]
which is zero.  Thus \(z\) has such a lift.

The group
\[
 \pi_{52\rho_K-d\lambda}^K(EK_+\wedge R_N)
\]
vanishes if \(52|K|-2d<0\).
Indeed, \(EK\) has free cells in nonnegative dimensions.
A free \(q\)-cell contributes an induced spectrum, whose
group in representation degree \(52\rho_K-d\lambda\) is
the underlying group of \(R_N\) in degree
\(52|K|-2d-q\).  It is zero because
\(\operatorname{Res}_eR_N=\MU^{\wedge |K|}\) is connective.
Induction over finite skeleta, followed by compactness of a
representation sphere, proves the vanishing.
Multiplication of the lifted class by \(a_\lambda^d\)
therefore gives
\begin{equation}\label{det:euler-times-norm}
 a_\lambda^d z=0
 \quad\text{when}\quad 52|K|-2d<0.
\end{equation}
The least such exponents are \(d=79\) for \(H\) and
\(d=235\) for \(G\).

\emph{The image of \(z\) is invertible.}
Let \(E^{\mathrm{Bor}}\) denote the cofree genuine
\(K\)-spectrum associated to the given action.
There is a map
\[
 R_N\longrightarrow E^{\mathrm{Bor}}
\]
obtained on underlying spectra with action by applying the
translated orientation in each factor and multiplying:
\[
 (a_\gamma)_{\gamma\in K}\longmapsto
          \prod_{\gamma\in K}\gamma f(a_\gamma).
\]
The coherent action and the \(E_\infty\)-multiplication
make this a map with coherent action; cofree adjunction
gives the genuine map.
Its underlying map is homotopy multiplicative, which is
all that will be needed.  In particular, an \(E_\infty\)
refinement of the orientation has not been assumed.

Let \(c_E\) be the underlying image of \(c\).
On the special fibre, \(\sigma-1\) has \(S\)-valuation \(3\),
so its first nonzero nonlinear strict coefficient occurs
at \(x^{27}\).  Hence \(c_E\) has unit residue and is a
homogeneous unit in \(E_*\).
The underlying image of \(z\) is the product of its
\(K\)-conjugates, and is again a unit.
Multiplication by \(z_E\) is thus an underlying equivalence
\[
 S^{52\rho_K}\wedge E^{\mathrm{Bor}}
       \longrightarrow E^{\mathrm{Bor}}.
\]
Both spectra are cofree, since a finite representation
sphere preserves cofreeness.  An underlying equivalence
between cofree spectra is a genuine equivalence.
Transport \eqref{det:euler-times-norm} along the map of
spectra and cancel multiplication by \(z_E\).
This proves the asserted vanishing of the actual Euler
powers.

\emph{The finite-page bound.}
The Tate \(E_2\)-unit is nonzero: the equivariant residue map
\(R\to k\) takes every additive trace to zero, since \(K\)
acts trivially on \(k\) and \(|K|=0\) there.  Cyclic
periodicity therefore makes every Euler-power symbol nonzero
on the ordinary \(E_2\)-page.
The skeletal filtration-two lift of \(a_\lambda\) has
symbol \eqref{det:euler-symbol}.
Multiplying these lifts with a cellular diagonal gives
a lift of \(a_\lambda^d\) in filtration \(2d\), with the
corresponding power as symbol.
It maps to zero in filtration zero by the result just proved.
Lemma~\ref{det:finite-boundary} therefore makes the symbol
an incoming boundary by length \(2d\).
Its representation internal degree is \(d(2-\lambda)\),
of underlying dimension zero.
An incoming \(d_r\) has source internal dimension \(1-r\),
so evenness forces \(r\) odd.
The bounds improve to \(2d-1\), namely \(157\) and \(469\).
After inverting the Euler symbol, a boundary on its power
is a boundary on the unit.  This proves the claim entirely
on finite pages.
\end{proof}
